% \iffalse meta-comment
%
% hit-final-abstract.dtx 是中英文摘要页
%
% \fi
%
% \section{中英文摘要页}
%
% \changes{v3.2a}{2026/09/28}{新增 \env{abstract} 环境：写在哪里就在哪里排一篇摘要（\option{lang}、\option{keywords}、
%   \option{shown}），与 iota-hit 的 \texttt{abstract()} 对应；v3.1e 的 \env{cabstract}、\env{eabstract} 只给老写法的
%   \cs{makecover} 收正文}
% \changes{v3.2a}{2026/09/28}{英文关键词标签只把“Keywords”加粗，后面跟细的全角冒号，照范例与 iota；博士后照模板仍是粗的“Keywords:”}
% \changes{v3.2a}{2026/09/27}{关键词换行后默认顶格，不再悬挂到标签之后：学校的 Word 模板与范例都是
%   普通段落；要悬挂写 \env{abstract}\oarg{keywords-hanging}（照 iota-hit 是这一篇摘要的参数）。“关键词：”的冒号
%   照 Word 占整格（装盒时 \pkg{luatexja} 只给半格）}
% \changes{v3.2a}{2026/09/27}{符号表与缩略语表的第一列按内容定宽（与 iota-hit 同一个算法：在 2.5\,cm
%   与各标签自己的宽里挑整张表最矮的那个，标签溢出的每行记一行的罚），\option{labelwidth} 写了才钉死；
%   合并成一页时两段取同一个宽}
% \changes{v3.2a}{2026/09/25}{摘要标题的 1.25 倍行距直接设在标题这一段上；博士后的关键词照模板，
%   不带冒号、后空一格，英文摘要单倍贴网格}
% \changes{v3.2a}{2026/08/22}{\cs{hit@make@abstract} 迁到 expl3（写进目录的 \texttt{\string~}
%   换成 \cs{nobreakspace}）；符号表提成 \env{symbols}，\env{denotation} 退回 v3.1e 原样
%   只供老文档；每篇摘要排完清一次缩写的「用过」状态，摘要里的缩写跟这篇摘要的语种排}
%
%    \begin{macrocode}
%<*final-abstract>
\bool_if:NT \g__hit_final_bool
  {
\box_new:N \hit@keywords@box
\NewDocumentCommand \hit@put@keywords { m m }
  {
    \group_begin:
      \docxlike_format_run_hbox_set:Nn \hit@keywords@box {#1}
      \noindent
      \bool_if:NT \l__hit_keywords_hanging_bool
        { \hangindent \wd \hit@keywords@box \hangafter 1 }
      \box \hit@keywords@box #2 \par
    \group_end:
  }

%    \end{macrocode}
%
% \changes{v3.0.0}{2020/05/21}{英文版目录之前有page行}
% \changes{v3.0.0}{2020/05/21}{博后关键字加粗}
%
% 两个摘要标题照指南排 1.25 倍行距，是标题这一段上的直接格式，排完就没了。英文摘要的标题
% 全是西文，标题自己会按西文行盒排（Word 按行里实际的字算行高）。
%
% 博士后的“关键词”照《博士后研究工作报告模板》：黑体加粗，不带冒号，后面敲两个半角空格。
% 标签装在盒里，盒里的空格不按贴着汉字各占半格排，所以直接空一格（两个半格）；盒里末字
% 后面也没有字距增量，一并补上。
%    \begin{macrocode}
\cs_new_protected:Npn \__hit_abstract_heading:n #1
  {
    \docxlike_format_with:nnn { Heading1 }
      { paragraph = { spacing = { line-spacing = multiple , at = 1.25 } } }
      {#1}
  }
\bool_new:N \g__hit_abstract_started_bool
\cs_new_protected:Npn \__hit_abstract_open:
  {
    \hit@clear@page
    \hit@layout@push@part:n { frontmatter }
    \bool_lazy_and:nnT { ! \g__hit_abstract_started_bool } { \g__hit_en_bool }
      {
        \addtocontents { toc } { \nobreakspace { } \hfill \textbf { \hit@term@table@of@contents@page@en } \par }
        \addtocontents { \ext@table }
          { { \noindent \sihao \textbf { \hit@term@table@of@contents@table@en } \nobreakspace { } \hfill \textbf { \hit@term@table@of@contents@page@en } } \par }
        \addtocontents { \ext@figure }
          { { \noindent \sihao \textbf { \hit@term@table@of@contents@figure@en } \nobreakspace { } \hfill \textbf { \hit@term@table@of@contents@page@en } } \par }
      }
  }
\bool_new:N \l__hit_abstract_restyle_bool
\cs_new_protected:Npn \__hit_abstract_numbering:
  {
    \bool_if:NF \g__hit_abstract_started_bool
      {
        \bool_gset_true:N \g__hit_abstract_started_bool
        \bool_set_true:N \l__hit_abstract_restyle_bool
        \pagestyle { hit@headings }
        \pagenumbering { Roman }
      }
  }
\cs_new_protected:Npn \__hit_abstract_restyle:
  { \pagestyle { hit@headings } }
%    \end{macrocode}
%
% 中文摘要排完先送出这一页再还版面：原先两篇摘要在同一次版面切换里排，中间那次换页落在
% 切换之内；英文摘要末尾不换页，由下一个部件送出，也与原先一样。
%    \begin{macrocode}
\cs_new_protected:Npn \__hit_abstract_zh:nn #1#2
  {
    \__hit_abstract_open:
    \__hit_abstract_heading:n
      {
        \hit@appendix@chapter * [ \hit@term@abstract@heading@zh ]
          { \hit@spread@by { \hit@term@abstract@heading@spread } { \hit@term@abstract@heading@zh } }
          [ \hit@term@abstract@in@english@toc@zh ]
      }
    \__hit_abstract_numbering:
    \tl_gset:Nn \g__hit_part_lang_tl { zh }
    #1
    % 正文与关键词之间按范例是一个空段(单独回车一行),\hitenter 模拟之
    % (1-8 页版 p4,36.8 = 19.4 + 17.25)
    \hitenter
    \hit@put@keywords
      {
        \bool_if:NT \g__hit_postdoc_bool { \bfseries } \heiti \hit@term@abstract@keywords@zh
        \bool_if:NT \g__hit_postdoc_bool
          {
            \skip_horizontal:n
              { 2 \g_docxlike_page_char_dim - \g_docxlike_body_size_dim }
          }
      }
      {#2}
    \hit@clear@page
    \tl_gclear:N \g__hit_part_lang_tl
    \__hit_abstract_reset:
    \hit@layout@pop@part:
  }
\cs_new_protected:Npn \__hit_abstract_en:nn #1#2
  {
    \__hit_abstract_open:
    \tl_gset:Nn \g__hit_part_lang_tl { en }
    \__hit_abstract_heading:n
      { \hit@appendix@chapter * { \hit@term@abstract@heading@en } [ \hit@term@abstract@in@english@toc@en ] }
    \__hit_abstract_numbering:
    % 英文摘要正文按范例排:Word 的 1.25 倍自动行距按行内字体自然行高算,
    % 纯西文段 TNR 12bp 实测行距 17.25;字符网格的行末余量 \rightskip 只属于
    % 汉字网格行,西文段排满真版心(例右缘 510.4);汉字行撑竿会把行距顶回
    % 网格,组内停用;范例导出勾了不断字,断词罚分同关
    \group_begin:
      % 按 abstract-en 目标铺版面(行距/版宽/首行基线/断词,见 hit-format)
      \docxlike_format_apply_par:n { abstract-en }
      #1
      \hitenter
      \hit@put@keywords
        {
          \bool_if:NTF \g__hit_postdoc_bool
            { \textbf { \hit@term@abstract@keywords@en : \enskip } }
            { \textbf { \hit@term@abstract@keywords@en } \hit@term@abstract@keywords@colon@en }
        }
        {#2}
    \group_end:
%    \end{macrocode}
%
% 摘要里出现过的缩写，正文第一次再出现时仍该全写一遍：读者未必先读摘要，
% 而且规范说的「文中第一次出现」指的是正文。两篇摘要之间也清一次，英文摘要
% 与中文摘要各自从头算，与 iota-hit 每篇摘要开头清一次同效。清不清由
% \option{abbreviation-reset-after-abstract} 定，默认清。
%    \begin{macrocode}
    \tl_gclear:N \g__hit_part_lang_tl
    \__hit_abstract_reset:
    \hit@layout@pop@part:
  }
%    \end{macrocode}
%
% \begin{environment}{abstract}
% \env{abstract}\oarg{选项} 写在哪里就在哪里排一篇摘要，与 iota-hit 的 \texttt{abstract()} 对应。
% 选项：\option{lang} 取 \texttt{zh}（默认）或 \texttt{en}；\option{keywords} 是这篇的关键词，逗号隔开，
% 不写就用元信息里的 \option{keywords-zh}、\option{keywords-en}；\option{shown} 同别的部件（摘要是
% 终稿专有的）。中英两篇各写一个环境，各自只排一次。环境体是整个抓下来再排的，里面不能放
% \cs{verb}。
%    \begin{macrocode}
\tl_new:N \l__hit_abstract_lang_tl
%    \end{macrocode}
%
% \option{keywords-hanging} 管这一篇的关键词折行后从哪起排：\texttt{false}（默认，\texttt{auto} 同义）顶格，与学校的
% Word 模板一致，那里关键词只是一个普通段落；\texttt{true} 悬挂到“关键词：”之后，各行的词左端对齐。
%    \begin{macrocode}
\bool_new:N \l__hit_keywords_hanging_bool
\keys_define:nn { hit / abstract }
  {
    keywords-hanging .choice: ,
    keywords-hanging / true .code:n = { \bool_set_true:N \l__hit_keywords_hanging_bool } ,
    keywords-hanging / false .code:n = { \bool_set_false:N \l__hit_keywords_hanging_bool } ,
    keywords-hanging / auto .code:n = { \bool_set_false:N \l__hit_keywords_hanging_bool } ,
    keywords-hanging .default:n = true ,
    lang .choices:nn = { zh , en }
      { \tl_set_eq:NN \l__hit_abstract_lang_tl \l_keys_choice_tl } ,
    lang .value_required:n = true ,
    keywords .code:n =
      { \hitsetup { info = { keywords- \l__hit_abstract_lang_tl = {#1} } } } ,
    shown .code:n = { \hit@part@set@shown:n {#1} } ,
  }
\NewDocumentEnvironment { abstract } { O{} +b }
  {
    \tl_set:Nn \l__hit_abstract_lang_tl { zh }
    \tl_set:Nn \l_hit_part_shown_tl { auto }
    \__hit_abstract_lang_first:n {#1}
    \hit@keys@set:nn { hit / abstract } {#1}
    \hit@part@if:nVTF { abstract } \l_hit_part_shown_tl
      {
        \exp_args:Ne \hit@part@once:nT { abstract- \l__hit_abstract_lang_tl }
          {
            \hit@matter@front@enter:
            \str_if_eq:VnTF \l__hit_abstract_lang_tl { en }
              { \__hit_abstract_en:nn {#2} { \hit@info@keywords@en } }
              { \__hit_abstract_zh:nn {#2} { \hit@info@keywords@zh } }
          }
      }
      { }
    \bool_if:NT \l__hit_abstract_restyle_bool
      { \group_insert_after:N \__hit_abstract_restyle: }
  }
  { }
%    \end{macrocode}
%
% \emph{页面样式要留到环境外面。} \cs{pagestyle} 是局部的，第一篇摘要把前置的页面样式换成带页眉页脚的那一份，
% 写在环境里的话环境一结束就退回去，后面的目录页没了页眉页脚。原先摘要在 \cs{makecover} 里、处在最外层，
% 没有这个问题；现在环境结束后再设一次。
%    \begin{macrocode}
%    \end{macrocode}
%
% \option{keywords} 要按这篇的语种设，而键是按写的次序处理的，\texttt{keywords} 写在
% \texttt{lang} 前面就会设到中文那一格。所以先单独把 \option{lang} 挑出来。
%    \begin{macrocode}
\cs_new_protected:Npn \__hit_abstract_lang_first:n #1
  {
    \clist_map_inline:nn {#1}
      {
        \__hit_abstract_lang_pick:w ##1 = \q_mark = \q_stop
      }
  }
\cs_new_protected:Npn \__hit_abstract_lang_pick:w #1 = #2 = #3 \q_stop
  {
    \str_if_eq:eeT { \tl_trim_spaces:n {#1} } { lang }
      { \tl_set:Ne \l__hit_abstract_lang_tl { \tl_trim_spaces:n {#2} } }
  }
%    \end{macrocode}
% \end{environment}
%
% 老写法：\cs{makecover} 排完封面扉页接着排两篇摘要，正文取自 \env{cabstract}、\env{eabstract}。
%    \begin{macrocode}
\cs_new_protected:Npn \hit@make@abstract
  {
    \__hit_abstract_zh:nn { \hit@abstract@zh } { \hit@info@keywords@zh }
    \__hit_abstract_en:nn { \hit@abstract@en } { \hit@info@keywords@en }
  }
\cs_new_protected:Npn \__hit_abstract_reset:
  {
    \bool_if:NT \g__hit_abbr_reset_bool
      { \cs_if_exist:NT \glsresetall { \glsresetall } }
  }
%    \end{macrocode}
%
% 符号表。
%
% 两个环境分工：\env{symbols} 是新写法，\env{denotation} 是 v3.1e 的老写法。
%
% \env{symbols} 本身就是一个 \pkg{enumitem} 的 \texttt{description} 列表，
% 一行一个符号 \cs{item}\oarg{符号}\meta{说明}。可选参数收键值：
% \option{labelwidth} 是第一列宽度（不写就按内容定，见下），\option{subheading} 给这一段
% 加个小标题——只在一页里排两段（符号加缩略语）时才用得着。不含等号的裸值当
% 宽度，\cs{begin}\marg{symbols}\oarg{3cm} 照样认。
%
% \emph{第一列多宽按内容定}，与 iota-hit 的 \texttt{auto-label-width} 同一个算法：候选宽度是
% 2.5\,cm（\pkg{thuthesis}、\pkg{bithesis} 的 \texttt{labelwidth}）与各个标签自己的宽里
% 大于它、不过半个版心的那些；每个候选把整张表试排一遍，总高最小的胜，标签比列宽的每行
% 另记一行的罚（第一列尽量不溢，除非为它加宽整列反倒让别的行折得更多），同高取宽的。
% 试排是整张表排进盒子量高深，不排出来。合并成一页时两段要同宽，说明列才对齐：两段各自
% 定下的宽取大的那个记进 \file{.aux}，下一遍两段都用它。
%
% 规范里允许把符号分成几张表、各有各的列，列表排不出来，所以还有
% \env{symbols*}：标题和表号照旧，正文由用户自己写。
%
% \env{denotation} 不带列表机制，也没有星号版：v3.1e 的它就是「排标题、重置表号、
% 正文随你写」，老文档里的内容是一张手写的 \env{table}。既然 \env{symbols} 全盘
% 接手了新写法，它退回原样即可，不必现代化，也不发废弃提示。
%    \begin{macrocode}
\RequirePackage{enumitem}
\newlist { hit@symbols } { description } { 1 }
\setlist [ hit@symbols ]
  {
    nosep ,
    font = \normalfont ,
    align = left ,
    leftmargin = ! ,% 下面三个长度之和
    labelindent = 0pt ,
    labelwidth = 2.5cm ,
    labelsep* = 0.5cm ,
    itemindent = 0pt ,
  }
\tl_new:N \l__hit_symbols_width_tl
\tl_new:N \l__hit_symbols_subheading_tl
\keys_define:nn { hit / symbols / opt }
  {
    labelwidth .tl_set:N = \l__hit_symbols_width_tl ,
    subheading .tl_set:N = \l__hit_symbols_subheading_tl ,
    shown .code:n = { \hit@part@set@shown:n {#1} } ,
  }
\cs_new_protected:Npn \__hit_symbols_opt:n #1
  {
    \tl_set:Nn \l__hit_symbols_width_tl { auto }
    \tl_clear:N \l__hit_symbols_subheading_tl
    \tl_set:Nn \l_hit_part_shown_tl { auto }
    \tl_if_in:nnTF {#1} { = }
      { \hit@keys@set:nn { hit / symbols / opt } {#1} }
      { \tl_if_blank:nF {#1} { \tl_set:Nn \l__hit_symbols_width_tl {#1} } }
  }
%    \end{macrocode}
%
% 每个部件各排各的章标题。只有合并成一页的那个部件例外：它自己排一个合并
% 标题，然后把闸门顶上，里头的 \env{symbols} 与 \cs{listofabbreviations*}
% 就都不再排章标题了。
%    \begin{macrocode}
\bool_new:N \g__hit_frontlist_open_bool
\tl_new:N \g__hit_frontlist_pending_tl
\cs_new_protected:Npn \hit@frontlist@open:nn #1#2
  {
    \bool_if:NF \g__hit_frontlist_open_bool
      {
        \hit@clear@page
%    \end{macrocode}
%
% 英文版排英文标题，中文版排中文标题、英文那个进英文目录。跟结论、致谢那几处
% 一样按语种分支——常量的 \texttt{-zh}／\texttt{-en} 后缀说的是「这个部件的
% 中文名与英文名」，不是「中文版用前者、英文版用后者」。
%    \begin{macrocode}
        \bool_if:NTF \g__hit_en_bool
          { \hit@appendix@chapter * {#2} [#2] }
          { \hit@appendix@chapter * {#1} [#2] }
      }
%    \end{macrocode}
%
% 段小标题:自己写的优先,没写就取合并部件预置的那一个,取完清掉。
%    \begin{macrocode}
    \tl_if_empty:NT \l__hit_symbols_subheading_tl
      {
        \tl_set_eq:NN \l__hit_symbols_subheading_tl
          \g__hit_frontlist_pending_tl
        \tl_gclear:N \g__hit_frontlist_pending_tl
      }
    \tl_if_empty:NF \l__hit_symbols_subheading_tl
      { \hit@frontlist@subheading:V \l__hit_symbols_subheading_tl }
  }
\cs_new_protected:Npn \hit@frontlist@subheading:n #1
  { \par \addvspace { \baselineskip } \noindent { \bfseries #1 } \par \nobreak }
\cs_generate_variant:Nn \hit@frontlist@subheading:n { V }
\cs_new_protected:Npn \__hit_symbols_open:
  {
    \hit@frontlist@open:nn
      { \hit@term@list@of@symbols@heading@zh } { \hit@term@list@of@symbols@heading@en }
    \setcounter { table } { 0 }
    \cs_set:Npn \thetable { \arabic { table } }% 表编号为 1
  }
\cs_new_protected:Npn \__hit_symbols_close:
  {
    \cs_set:Npn \thetable { \arabic { chapter } - \arabic { table } }% 表编号为 7-1
    \setcounter { table } { 0 }
  }
\cs_new_protected:Npn \__hit_symbols_if:TF #1#2
  {
    \bool_if:NTF \g__hit_frontlist_open_bool
      {#1}
      {
        \hit@part@if:nVTF { list-of-symbols } \l_hit_part_shown_tl
          {
            \hit@part@done:nTF { list-of-symbols }
              { \msg_warning:nnn { hithesis } { part-twice } { list-of-symbols } #2 }
              {
                \hit@part@once:nT { list-of-symbols } { }
                \hit@matter@front@enter:
                \hit@nomenclature@check:
                #1
              }
          }
          {#2}
      }
  }
%    \end{macrocode}
%
% \begin{macro}{\hit@nomenclature@check:}
% 符号与缩略语既排了单独的表（\cs{listofsymbols}、\env{symbols}、\cs{listofabbreviations}），又排了合页
% （\cs{printnomenclature}），两边都照排，同一批条目出现两遍，提示一次。iota-hit 这里直接报错停下；
% 这边照 \hologo{LaTeX} 的习惯让编译走完。
%    \begin{macrocode}
\cs_new_protected:Npn \hit@nomenclature@check:
  {
    \bool_lazy_and:nnT
      { \hit@part@done_p:n { nomenclature } }
      {
        \bool_lazy_or_p:nn
          { \hit@part@done_p:n { list-of-symbols } }
          { \hit@part@done_p:n { list-of-abbreviations } }
      }
      { \__hit_warn_once:nnn { nomenclature-twice } { nomenclature-twice } { } }
  }
%    \end{macrocode}
% \end{macro}
%
%    \begin{macrocode}
\NewDocumentEnvironment { symbols } { O{} +b }
  {
    \__hit_symbols_opt:n {#1}
    \__hit_symbols_if:TF
      { \__hit_symbols_open: \hit@symbols@list:n {#2} }
      { }
  }
  { \__hit_symbols_close: }
%    \end{macrocode}
%
% \begin{macro}{\hit@symbols@list:n}
% \cs{hit@symbols@list:n}\marg{条目} 把一串 \cs{item} 排成两列表，第一列宽度是
% \texttt{auto} 时先按上面的算法定下来。缩略语表（\file{src/flow/hit-final-glossary.dtx}）
% 也调它。
%    \begin{macrocode}
\dim_const:Nn \c__hit_symbols_width_min_dim { 2.5cm }
\seq_new:N \g__hit_symbols_label_seq
\seq_new:N \l__hit_symbols_candidate_seq
\tl_new:N \l__hit_symbols_body_tl
\box_new:N \l__hit_symbols_box
\dim_new:N \l__hit_symbols_best_dim
\dim_new:N \l__hit_symbols_cost_dim
\dim_new:N \l__hit_symbols_best_cost_dim
\dim_new:N \g__hit_symbols_shared_dim
\dim_new:N \g__hit_symbols_shared_prev_dim
\cs_new_protected:Npn \hit@symbols@list:n #1
  {
    \tl_set:Nn \l__hit_symbols_body_tl {#1}
    \__hit_symbols_width_fix:
    \begin { hit@symbols } [ labelwidth = \l__hit_symbols_width_tl ]
      \l__hit_symbols_body_tl
    \end { hit@symbols }
  }
\cs_new_protected:Npn \__hit_symbols_width_fix:
  {
    \str_if_eq:VnT \l__hit_symbols_width_tl { auto }
      {
        \__hit_symbols_auto:
        \bool_if:NT \g__hit_frontlist_open_bool
          {
            \dim_gset:Nn \g__hit_symbols_shared_dim
              { \dim_max:nn { \g__hit_symbols_shared_dim } { \l__hit_symbols_best_dim } }
            \dim_set:Nn \l__hit_symbols_best_dim
              { \dim_max:nn { \l__hit_symbols_best_dim } { \g__hit_symbols_shared_prev_dim } }
          }
        \tl_set:Nx \l__hit_symbols_width_tl { \dim_use:N \l__hit_symbols_best_dim }
      }
  }
\cs_new_protected:Npn \__hit_symbols_measure:n #1
  {
    \vbox_set:Nn \l__hit_symbols_box
      {
        \begin { hit@symbols } [ labelwidth = #1 ]
          \l__hit_symbols_body_tl
        \end { hit@symbols }
      }
  }
%    \end{macrocode}
%
% 先试排一遍收各个标签的宽：列表里的 \cs{item} 换成一个先量标签再交还原来那个的版本。
% 标签的字与列表里排的一样是 \cs{normalfont}。
%    \begin{macrocode}
\cs_new_protected:Npn \__hit_symbols_auto:
  {
    \seq_gclear:N \g__hit_symbols_label_seq
    \vbox_set:Nn \l__hit_symbols_box
      {
        \begin { hit@symbols } [ labelwidth = \c__hit_symbols_width_min_dim ]
          \DeclareCommandCopy \__hit_symbols_item:w \item
          \RenewDocumentCommand \item { o }
            {
              \IfValueTF {##1}
                {
                  \hbox_set:Nn \l_tmpa_box { \normalfont ##1 }
                  \seq_gput_right:Nx \g__hit_symbols_label_seq
                    { \dim_use:N \box_wd:N \l_tmpa_box }
                  \__hit_symbols_item:w [ ##1 ]
                }
                { \__hit_symbols_item:w }
            }
          \l__hit_symbols_body_tl
        \end { hit@symbols }
      }
    \seq_clear:N \l__hit_symbols_candidate_seq
    \seq_put_right:Nx \l__hit_symbols_candidate_seq
      { \dim_use:N \c__hit_symbols_width_min_dim }
    \seq_map_inline:Nn \g__hit_symbols_label_seq
      {
        \bool_lazy_and:nnT
          { \dim_compare_p:nNn {##1} > { \c__hit_symbols_width_min_dim } }
          { ! \dim_compare_p:nNn {##1} > { \linewidth / 2 } }
          { \seq_put_right:Nn \l__hit_symbols_candidate_seq {##1} }
      }
    \seq_remove_duplicates:N \l__hit_symbols_candidate_seq
    \seq_sort:Nn \l__hit_symbols_candidate_seq
      {
        \dim_compare:nNnTF {##1} > {##2}
          { \sort_return_swapped: } { \sort_return_same: }
      }
    \dim_set_eq:NN \l__hit_symbols_best_dim \c__hit_symbols_width_min_dim
    \dim_set_eq:NN \l__hit_symbols_best_cost_dim \c_max_dim
    \seq_map_inline:Nn \l__hit_symbols_candidate_seq
      {
        \__hit_symbols_measure:n {##1}
        \dim_set:Nn \l__hit_symbols_cost_dim
          { \box_ht:N \l__hit_symbols_box + \box_dp:N \l__hit_symbols_box }
        \seq_map_inline:Nn \g__hit_symbols_label_seq
          {
            \dim_compare:nNnT {####1} > {##1}
              { \dim_add:Nn \l__hit_symbols_cost_dim { \baselineskip } }
          }
        \dim_compare:nNnF { \l__hit_symbols_cost_dim } > { \l__hit_symbols_best_cost_dim }
          {
            \dim_set:Nn \l__hit_symbols_best_dim {##1}
            \dim_set_eq:NN \l__hit_symbols_best_cost_dim \l__hit_symbols_cost_dim
          }
      }
  }
%    \end{macrocode}
%
% 合并页两段共用的宽：本遍取的大者文档结束时写进 \file{.aux}，下一遍读回来。
%    \begin{macrocode}
\cs_new_protected:Npn \hit@symbols@shared #1
  { \dim_gset:Nn \g__hit_symbols_shared_prev_dim {#1} }
\AddToHook { enddocument }
  {
    \dim_compare:nNnT { \g__hit_symbols_shared_dim } > { \c_zero_dim }
      {
        \iow_now:Nx \@auxout
          {
            \token_to_str:N \hit@symbols@shared
            { \dim_use:N \g__hit_symbols_shared_dim }
          }
      }
  }
%    \end{macrocode}
% \end{macro}
%
%    \begin{macrocode}
\bool_new:N \l__hit_symbols_shown_bool
\NewDocumentEnvironment { symbols* } { O{} }
  {
    \__hit_symbols_opt:n {#1}
    \__hit_symbols_if:TF
      { \bool_set_true:N \l__hit_symbols_shown_bool \__hit_symbols_open: }
      { \bool_set_false:N \l__hit_symbols_shown_bool \hit@part@skip: }
  }
  { \bool_if:NT \l__hit_symbols_shown_bool { \__hit_symbols_close: } }
%    \end{macrocode}
%
% \begin{macro}{\listofsymbols}
% 按 \option{symbols} 族的词条排符号表，与 iota-hit 的 \texttt{list-of-symbols()} 对应。中文版取词条的
% \option{zh}，英文版取 \option{en}。有组的先按全部词条定第一列的宽，再一组一段：组名当小标题，
% 各段第一列一样宽。\cs{hit@symbols@setup:} 也给合页用。
%    \begin{macrocode}
\NewDocumentCommand \listofsymbols { O{} }
  {
    \__hit_symbols_opt:n {#1}
    \__hit_symbols_if:TF
      { \__hit_symbols_open: \hit@symbols@setup: \__hit_symbols_close: }
      { }
  }
\tl_new:N \l__hit_symbols_all_tl
\tl_new:N \l__hit_symbols_chunk_tl
\tl_new:N \l__hit_symbols_chunk_name_tl
\seq_new:N \l__hit_symbols_chunks_seq
\cs_new_protected:Npn \hit@symbols@setup:
  {
    \tl_clear:N \l__hit_symbols_all_tl
    \tl_clear:N \l__hit_symbols_chunk_tl
    \tl_clear:N \l__hit_symbols_chunk_name_tl
    \seq_clear:N \l__hit_symbols_chunks_seq
    \seq_map_inline:Nn \g_hit_symbols_seq { \__hit_symbols_collect:nnnn ##1 }
    \__hit_symbols_chunk_flush:
    \seq_if_empty:NTF \l__hit_symbols_chunks_seq
      { \msg_warning:nn { hithesis } { symbols-empty } }
      {
        \tl_set_eq:NN \l__hit_symbols_body_tl \l__hit_symbols_all_tl
        \__hit_symbols_width_fix:
        \seq_map_inline:Nn \l__hit_symbols_chunks_seq { \__hit_symbols_chunk:nn ##1 }
      }
  }
\cs_new_protected:Npn \__hit_symbols_collect:nnnn #1#2#3#4
  {
    \tl_if_eq:nnTF {#2} { \q_mark }
      {
        \__hit_symbols_chunk_flush:
        \tl_set:Nn \l__hit_symbols_chunk_name_tl {#1}
      }
      {
        \tl_set:Nx \l_tmpa_tl
          {
            \exp_not:N \item [ { \exp_not:n {#2} } ]
            \bool_if:NTF \g__hit_en_bool { \exp_not:n {#4} } { \exp_not:n {#3} }
          }
        \tl_put_right:NV \l__hit_symbols_chunk_tl \l_tmpa_tl
        \tl_put_right:NV \l__hit_symbols_all_tl \l_tmpa_tl
      }
  }
\cs_new_protected:Npn \__hit_symbols_chunk_flush:
  {
    \bool_lazy_and:nnF
      { \tl_if_empty_p:N \l__hit_symbols_chunk_tl }
      { \tl_if_empty_p:N \l__hit_symbols_chunk_name_tl }
      {
        \seq_put_right:Nx \l__hit_symbols_chunks_seq
          {
            { \exp_not:V \l__hit_symbols_chunk_name_tl }
            { \exp_not:V \l__hit_symbols_chunk_tl }
          }
      }
    \tl_clear:N \l__hit_symbols_chunk_tl
    \tl_clear:N \l__hit_symbols_chunk_name_tl
  }
\cs_new_protected:Npn \__hit_symbols_chunk:nn #1#2
  {
    \tl_if_empty:nF {#1} { \hit@frontlist@subheading:n {#1} }
    \tl_if_empty:nF {#2} { \hit@symbols@list:n {#2} }
  }
%    \end{macrocode}
% \end{macro}
%
% v3.1e 的 \env{denotation}：排标题、把表号改成不带章号的形式，正文随用户写。
%    \begin{macrocode}
\cs_new_protected:Npn \__hit_denotation_open:
  {
    \tl_set:Nn \l_hit_part_shown_tl { auto }
    \hit@part@if:nVTF { list-of-symbols } \l_hit_part_shown_tl
      {
        \bool_set_true:N \l__hit_symbols_shown_bool
        \hit@matter@front@enter:
        \hit@clear@page
        \hit@appendix@chapter *
          { \hit@term@list@of@symbols@heading@zh } [ \hit@term@list@of@symbols@heading@en ]
        \setcounter { table } { 0 }
        \cs_set:Npn \thetable { \arabic { table } }
      }
      { \bool_set_false:N \l__hit_symbols_shown_bool \hit@part@skip: }
  }
\cs_new_protected:Npn \__hit_denotation_close:
  {
    \bool_if:NT \l__hit_symbols_shown_bool
      {
        \cs_set:Npn \thetable { \arabic { chapter } - \arabic { table } }
        \setcounter { table } { 0 }
      }
  }
\NewDocumentEnvironment { denotation } { O{2.5cm} }
  { \__hit_denotation_open: }
  { \__hit_denotation_close: }
%    \end{macrocode}
%
% 带星号的不开列表，正文完全交给用户。可选参数收下不用，这样 v3.1e 里
% \cs{begin}\marg{denotation}\oarg{2.5cm} 的写法加个星号就能直接跑。
%    \begin{macrocode}
\NewDocumentEnvironment { denotation* } { O{2.5cm} }
  { \__hit_denotation_open: }
  { \__hit_denotation_close: }
%    \end{macrocode}
%
%    \begin{macrocode}
  }
%</final-abstract>
%    \end{macrocode}
%
% \Finale
